\section{Fourier 变换 \& 空间解析几何}

\subsection{Fourier 级数的复数形式}

定义复数 Fourier 系数：
$$
C_n = \frac{1}{2l} \int_{-l}^{l} f(x) \E^{-\frac{\I n \pi x}{l}} \dd{x}, \quad n = 0, \pm 1, \pm 2, \dots
$$
那么 Fourier 级数可以写为：
$$
f(x) = \sum_{n = -\infty}^{\infty} C_n \E^{\frac{\I n \pi x}{l}}
$$

\subsection{非周期函数的 Fourier 变换}

对于在 $\mathbb{R}$ 上可积的非周期函数 $f(x)$，我们做如下操作：

\begin{enumerate}
    \item 提取出其在 $[-l, l]$ 的部分，以 $2l$ 为周期延拓出周期函数 $\tilde{f}_l(x)$
    \item 对 $\tilde{f}_l(x)$ 作复数形式的 Fourier 级数展开：
    $$
    \tilde{f}_l(x) = \sum_{n = -\infty}^{+\infty} C_n(l) \E^{\I \omega_n x}
    $$
    其中：
    $$
    \omega_n = \frac{n \pi}{l}, \quad C_n(l) = \frac{1}{2l} \int_{-l}^{l} f(x) \E^{-\I \omega_n x} \dd{x}
    $$
    \item 令 $l \to +\infty$，那么得到：
    $$
    f(x) = \lim_{l \to +\infty} \tilde{f}_l(x) = \lim_{\Delta \omega \to 0} \frac{1}{2\pi} \sum_{-\infty}^{+\infty} \ab(\int_{-l}^{l} \tilde{f}_l(x) \E^{-\I \omega_n x} \dd{x}) \E^{\I \omega_n x} \Delta \omega
    $$
    记
    $$
    \varphi_l(\omega) = \int_{-l}^{l} \tilde{f}_l(x) \E^{-\I \omega x} \dd{x}
    $$
    那么
    $$
    f(x) = \lim_{\Delta \omega \to 0} \frac{1}{2\pi} \sum_{n = -\infty}^{+\infty} \varphi_l(\omega_n) \E^{\I \omega x} \Delta \omega = \frac{1}{2\pi} \int_{-\infty}^{+\infty} \varphi(\omega) \E^{\I \omega x} \dd{\omega}
    $$
    其中
    $$
    \varphi(\omega) = \lim_{l \to \infty} \int_{-l}^{l} \tilde{f}_l(x) \E^{-\I \omega x} \dd{x} = \int_{-\infty}^{+\infty} f(x) \E^{-\I \omega x} \dd{x}
    $$
\end{enumerate}

于是，我们定义：

\begin{definition}[Fourier 变换]
    \textbf{Fourier 正变换}为：
    $$
    \hat{f}(\omega) = \mathcal{F}\ab[f](\omega) = \frac{1}{\sqrt{2\pi}} \int_{-\infty}^{+\infty} f(x) \E^{-\I \omega x} \dd{x}
    $$
    \textbf{Fourier 逆变换}为：
    $$
    f(x) = \mathcal{F}^{-1}\ab[\hat{f}](x) = \frac{1}{\sqrt{2\pi}} \int_{-\infty}^{+\infty} \hat{f}(\omega) \E^{\I \omega x} \dd{\omega}
    $$
\end{definition}

\subsection{多元函数的 Fourier 变换}

\begin{definition}[多元 Fourier 变换]
    记 $\vect{\omega} = (\omega_1, \dots, \omega_n), \vect{x} = (x_1, \dots, x_n)$，若 $n$ 元函数 $f \in \mathbb{L}^{1}$，那么称 Fourier 正变换为：
    $$
    \hat{f}(\vect{\omega}) = \mathcal{F}\ab[f](\vect{\omega}) = \frac{1}{(2\pi)^{n/2}} \int_{\mathbb{R}^n} f(\vect{x}) \E^{-\I \vect{\omega} \cdot \vect{x}} \dd{\vect{x}}
    $$
    Fourier 逆变换为：
    $$
    f(\vect{x}) = \mathcal{F}^{-1}\ab[\hat{f}](\vect{x}) = \frac{1}{(2\pi)^{n/2}} \int_{\mathbb{R}^n} \hat{f}(\vect{\omega}) \E^{\I \vect{\omega} \cdot \vect{x}} \dd{\vect{\omega}}
    $$
\end{definition}

\subsection{Fourier 变换的性质}

\subsubsection{线性性}

对于任意的复数 $\alpha, \beta \in \mathbb{C}$，都有：
$$
\begin{gathered}
    \mathcal{F}\ab[\alpha f(x) + \beta g(x)](\omega) = \alpha \mathcal{F}\ab[f](\omega) + \beta \mathcal{F}\ab[g](\omega), \\
    \mathcal{F}^{-1}\ab[\alpha \hat{f}(\omega) + \beta \hat{g}(\omega)](x) = \alpha \mathcal{F}^{-1}\ab[\hat{f}](x) + \beta \mathcal{F}^{-1}\ab[\hat{g}](x)
\end{gathered}
$$

\subsubsection{互逆性}

显然，Fourier 正变换和逆变换满足如下的\textbf{互逆性}：
$$
\begin{gathered}
    F^{-1}\ab[\mathcal{F}\ab[f]](x) = f(x), \\
    \mathcal{F}\ab[F^{-1}\ab[\hat{f}]](\omega) = \hat{f}(\omega)
\end{gathered}
$$

\subsubsection{微分性质}

如果 $\abs{x} \to \infty$ 时有 $f(x) \to 0,\ f^{(n - 1)}(x) \to 0$，那么
$$
\mathcal{F}\ab[f^{(k)}](\omega) = (\I \omega)^k \mathcal{F}\ab[f](\omega), \quad k = 1, 2, \dots, n
$$

\subsection{Fourier 变换的应用}

\subsubsection{求解常微分方程组}

使用 Fourier 变换，微分运算可以转换为代数运算，从而简化方程的求解。

\subsection{空间平面及其方程}

\begin{definition}[法向量]
    若一个非零向量 $\vect{n}$ 垂直于平面 $\alpha$，那么称 $\vect{n}$ 为平面 $\alpha$ 的一个\textbf{法向量}。
\end{definition}

\subsubsection{点法式方程}

对于平面 $\alpha$，点 $M(x, y, z) \in \alpha$，若有点 $M_0(x_0, y_0, z_0) \in \alpha$ 与法向量 $\vect{n} = (n_x, n_y, n_z)$，那么就有：
$$
\overrightarrow{M_0 M} \perp \vect{n} \iff \overrightarrow{M_0 M} \cdot \vect{n} = 0
$$
于是得到方程：
$$
n_x (x - x_0) + n_y (y - y_0) + n_z (z - z_0) = 0
$$
这就是空间平面的\textbf{点法式方程}。

\subsubsection{一般方程}

对于方程
$$
A x + B y + C z + D = 0 \quad \ab(A^2 + B^2 + C^2 \ne 0)
$$
其称为空间平面的\textbf{一般方程}。不难发现，其法向量为
$$
\vect{n} = (A, B, C)
$$
